\documentclass[../main]{subfiles}
\begin{document}
\ifSubfilesClassLoaded{\setcounter{page}{50}}{}
\section{Коммунизм и семья}
\label{sec:09_kommunizm_semya}

Коммунизм и семья «Согласно материалистическому пониманию, определяющим моментом в истории является в конечном счете производство и воспроизводство непосредственной жизни. Но само оно, опять-таки, бывает двоякого рода. С одной стороны - производство средств к жизни: предметов питания, одежды, жилища и необходимых для этого орудий; с другой - производство самого человека, продолжение рода. Общественные порядки, при которых живут люди определенной исторической эпохи и определенной страны, обусловливаются обоими видами производства: ступенью развития, с одной стороны - труда, с другой - семьи», - писал Фридрих Энгельс в работе «Происхождение семьи, частной собственности и государства».

В предыдущих главах мы говорили о том, что при коммунизме производство жизненно необходимых предметов автоматизировано и подчинено производству человеческой жизни во всей её полноте. Общество сознательно творит своё собственное развитие, свои собственные формы коллективности, расширяя содержание свободной деятельности. Семья, как хозяйственная ячейка, основанная на частном характере присвоения продуктов труда, утрачивает своё значение и исчезает в переходный период. «Частное домашнее хозяйство превратится в общественную отрасль труда», - писал Энгельс. Выше мы видели насколько это экономит общественное время, освобождая людей от тягот быта. Однако обобществление касается не только потребления.

Главный переворот в области семейных отношений состоит в общественном воспитании детей. Этот переворот давно назревает, и движение в этом направлении есть уже при капитализме. Домашнее хозяйство начало «обобществляться», быт давно во многом стал отраслью общественного хозяйства. Мы готовим еду в основном из полуфабрикатов, покупаем одежду в магазине, стираем её в прачечной, хотя совсем недавно это было женским делом, совершающимся в домашних условиях. Современному мужчине тоже не нужно носить воду, колоть дрова, и строить дом, в котором живёт семья. Всё это давно является общественным высокотехнологичным делом.

Но вот воспитание детей все никак не превратится в сферу производства, хотя общественные заведения для этого появились давно, а активно развиваются уже не меньше двух столетий. И всегда те, кто побогаче, отдавал детей именно туда, нанимал частных преподавателей, а сам лишь получал удовольствие от семейного уюта. Также и в Советском Союзе – одним из первых достижений, о которых с такой гордостью говорили социалистические трудящиеся, это организация яслей и детских садов при фабриках. То есть потребность отдать воспитание детей тем, кто это умеет, выиграть свободное время для того, что считаешь действительно нужным и важным (для кого-то, кстати, это и воспитание детей, почему нет) – в этом едины и бедные, и богатые, с некоторыми лишь различиями мещанско-поэтического характера. Капитализм давным-давно рушит традиционную семью, и этот процесс не имеет моральных корней, но только экономические. В настоящий момент пропаганда семейных ценностей обусловлена с одной стороны, достижениями социалистических революций, запретивших эксплуатацию детского труда, а с другой – потребностью капитала в рынке детских товаров и образования. Но действительной исторической, общественной необходимости      в   сохранении   семьи   в   нынешней   форме хозяйственной ячейки нет.

Итак, с переходом общества к коммунистическим отношениям, сфера воспитания детей станет сферой общественного производства. Воспитание становится подлинным социальным творчеством - свободной деятельности по меркам истины, добра и красоты. Это во много раз повышает ответственность воспитателей и перед детьми, и перед обществом в целом. Как писал когда-то А.С. Макаренко: «Никаких детских катастроф, никаких неудач, никаких процентов брака, даже выраженных сотыми единицы, у нас быть не должно!» В ряде работ он сравнивал воспитание с другими отраслями и показывал, насколько неудача здесь опаснее, чем где бы то ни было. Говоря о семье, Антон Семенович замечал, что задача родителей – вырастить достойных членов общества, за «качество поставляемого материала» они несут перед обществом ответственность, получая взамен теплоту семейных отношений. Но ведь далеко не все родители умеют воспитывать детей. Вдобавок далеко не у всех детей есть родители или просто полноценная семья уже сейчас. В этом заключается необходимость в сознательном, качественном обобществлении воспитания, в создании иных форм «семьи».

Коммунистическая «семья» - это коллективная форма, в которой каждый причастен к производству общественной жизни в общественных же масштабах. Из этого ясно, что размеры «семьи» могут варьироваться, и уж точно она не может замыкаться в современных рамках. Размеры зависит от многих конкретно- исторических условий общественного производства, а потому будут возможны различные варианты, как в рамках одной эпохи в развитии коммунистического общества, так и в его истории. Поскольку такая «семья» как первичный коллектив становящегося человека будет основана на общественном характере деятельности и непосредственно определяться им, это значит, что принцип создания таких коллективов тоже будет непосредственно-общественным - качественно иным, чем сейчас. Речь идёт об уничтожении современного брака как формы, тем более, что половые отношения и половая любовь и брачно- хозяйственные отношения — это далеко не одно и то же. Поэтому одним из важнейших и непреходящих практических открытий А.С. Макаренко было то, что новая «семья» должна напрямую включать каждого человека в общественное производство, как активного члена общественного трудового коллектива, а не как потребителя общественных благ ил как «объект» воспитания.

Отслеживая быстрые, порой незаметные невооружённому глазу изменения в этой сфере в первые же годы Великой Октябрьской социалистической революции, Александра Коллонтай в своей работе «Дорогу крылатому Эросу!» попыталась увидеть и спрогнозировать тот тип личных, а точнее отношений половой любви, который утвердится при полной победе общественной собственности над частной, при развитии новых общественных институтов и при внедрении элементов инфраструктуры, способной разгрузить материнство, всё более превращая его в отрасль общественного производства.

Новый быт, там, где он был налажен (не только в СССР), действительно раскрепостил женщин в невообразимых для всего остального мира масштабах, дал время не только на саму любовь и «ухаживания», но и на переосмысление любви. Такие новые, уже частично коммунистические отношения описаны и в художественной литературе времён индустриализации, как в прозе самой Александры Коллонтай («Василиса Малыгина», «Большая любовь»), так, например, у Ильи Эренбурга в романе «День второй».

Утверждение коммуны без регистрации брака (брак без регистрации стал в 1920-1930-х нормой, регистрация - наоборот, признаком недоверия в паре) там, где до этого царствовала семья патриархального типа, конечно, было сильным опережением событий. Но это было революцией в самовосприятии любящих, влюблённых, семейных, разведённых. Не брак как жизнь для двоих, ради двоих, а семейная жизнь, непосредственно определяемая целями бурно растущего и развивающегося общества, ближайшего (трудового) коллектива – вот реальная повестка любви, выражающая тенденцию к коммунизму.

Семья как зона любви и связанной с ней ответственности в капиталистическом обществе – «зона повышенного риска». Супруги обречены в стенах своего жилья переживать всё то, что запрограммировано отношениями частной собственности, вопросами достатка, заработка, трат и прочей прозы жизни. Но коммунизм и есть устранение таких проблем.

Динамика отмирания и прежней семьи, и государства как такового, сложна. Однако, полумеры социализма тут бессильны: из коммун-общежитий студенты, молодые специалисты, всё же оказывались в квартирах, в семьях, в старых рамках, но с новым соцобеспечением. Мотивы перерастать прежние ячейки «сот» - растерялись в обществе, где денежно-товарные отношения опять отвоевывали себе господствующее положение. За рывком в коммуну последовал закономерный и тоже будто бы временный возврат к патриархальности. Снова забота только родителей, а не всего коллектива или общества, выходила в воспитании на первый план. А ведь намечались не только «сыны полка» (война вносила коррективы), но и дети коммун.

В этих «детях коммун» и воплощался один из важнейших моментов для становления коммунизма. При переходе к коммунизму дети должны перестать делиться на своих и чужих. Не существует чужих детей (то есть детей, о которых должен заботится кто-то другой, но не я) в громадном коллективе, где сняты проблемы «счётчика» личного трудового вклада.

Оказываясь изначально в коллективе, где много мужчин и женщин, дети, различая отца и мать, не будут видеть «край и конец света» там, где мама и папа заканчиваются и начинаются «соседи». Через естественную в коммуне социализацию, через повседневное ощущение на себе заботы коллектива, дети перестанут «доигрывать в песочнице частной собственности». Уход от старого брака позволит избежать ряда пережитков, воспроизводимых в психике, в сознании отцов и матерей, а не только детей. Смена парности в большом и дружном коллективе перестаёт быть чем-то травмирующим, если между членами пары уже нет любви. Вообще при отсутствии общественных экономических барьеров, всё сугубо личное (в том числе и интимное) только и может стать подлинно личным, но не противостоящим общественному, освободившись от хозяйственных, правовых и моральных оков.

Новая «семья» сама должна стать непосредственно-коллективной формой коммунистической деятельности по созданию нового человека, непосредственно связывающая всех её членов, включая детей с обществом как деятельных субъектов.
\end{document}
